% !TeX program = xelatex
% 本例的寫本與相關讀法純屬虛構，僅展示書稿版式。
\documentclass[lang=zh]{impebook}
\UseTemplateSet{features={image,hyperlinks}}
\title{寫本文獻比較研究示例}
\author{作者姓名}
\date{\today}
\begin{document}
\frontmatter
\maketitle
\chapter*{前言}
本書是一份用來展示學術書稿排版的小型範例。書中甲、乙、
丙三件寫本及其文句均屬假設材料，目的在於讓讀者觀察
扉頁、前言、目錄、章首頁、頁眉、腳註、正文段落與
交叉引用的實際效果。它並不企圖對某一部古籍提出
新的研究結論，更不能取代對實際文獻的整理與考證。

長篇書稿與簡短的測試文件不同。幾個孤立的句子無法
充分展示行距、段落縮排和長篇正文的閱讀節奏，
也難以看清章節之間如何分頁。因此，本例圍繞一個
連續的問題展開：研究者如何由不同傳本出發，
逐步形成一份可受檢驗的校訂文。各章篇幅適中，
適合初步比較版面效果。

本文盡量採用 \LaTeX{} 的常規命令。初學者可以
直接替換文字與章節，不必先掌握大量特殊接口；
需要書目、索引、圖片或古文字時，再依照文獻
性質增加相應功能即可。

\tableofcontents
\mainmatter

\chapter{研究對象與問題}\label{chap:problem}
\section{何謂文獻傳本}
我們所接觸的古代文本，通常寄寓於具體的物質
載體之中。簡牘、帛書、紙本、碑刻或傳世印本
可以保存相近的內容，卻很少在所有字句上完全
一致。部分差異來自材料殘損，部分差異與抄寫、
改編及讀者需求有關。研究者若只選出其中一個
版本，而不交代其他材料，便難以說明文本
最終形態是如何形成的。

本例假設甲、乙、丙三件寫本保存同一短句。
甲本較早，文字簡略；乙本較晚，語句較長；
丙本關鍵位置殘缺。甲乙丙只是便於討論的
代號，不代表證據強弱的排序。晚出的傳本
有時可能保存古老的讀法，較早的寫本也
可能有獨特的訛字或刪省。

在正式研究中，首先需要逐件交代文獻的
收藏情況、形制、字體、書寫層次與年代
依據。若材料曾經整理、拼綴或重新編號，
也應說明相應的研究史。這類資訊雖然
繁瑣，卻直接決定後來的校訂能否由他人
覆核。

\section{從字形到錄文}
錄文不是對材料的機械複製。研究者必須
辨認筆畫、判斷文字、處理斷簡，並以
既定符號記錄殘損和補寫。若將推測所得
的文字不加說明地填入原文，校訂階段
便容易形成循環論證。因此，最好區分
忠實記錄字形的材料錄文與便於理解的
通讀文本，並說明兩者之間的轉換原則。

例如，丙本中段有一處缺失，不能僅因
甲本在相應位置寫了某字，就將該字
當成丙本的既有內容。缺失是材料的現狀，
補字則是編者的判斷。二者都值得記錄，
卻屬於不同層次。\footnote{實際的文獻
整理應說明闕文、重文、合文、衍文及補字
等記號的具體使用方法，並盡可能提供圖版。}

\section{年代與可信度}
寫本的年代是重要資料，卻不能直接等同
於文字的可靠程度。甲本即使確實早於乙本，
也不意味着其所有字句都更接近某個
假設的祖本。要考察各本之間的關係，
還需要比較共同訛誤、相同增省與特殊
用語，並考慮是否存在亡佚的中間傳本。

下一章將在這些限制之下，對三件寫本
進行比較。先描述文獻，再分析異文，
有助於避免研究者在整理材料以前
便作出喜好的選擇。當具體證據不充分時，
暫時保留多種可能性也是合理的學術結論。

\chapter{異文的比較}\label{chap:variants}
\section{增字與闕文}
假設甲本作「謹守其書」，乙本作
「敬謹守其書」，丙本僅保存「謹……其書」。
從表面看，乙本比甲本多出一個修飾字，
而丙本無法判斷這個位置究竟曾經寫有
甚麼。乙本的「敬」字或許出於抄手
對文句的補充，也可能是甲本傳抄時
省略了本有之字。

丙本不應被隨意納入支持甲本的證據。
它只是與較短的讀法不相衝突，卻同樣
可以容納較長的形式。將闕文當成某一
特定異文的證明，是文獻比勘中常見
而需要警惕的推理錯誤。因此，記錄
一項尚未解決的疑問，有時比作出斷然
的校訂更能反映材料的真實狀態。

\section{比較所依據的標準}
研究者可以藉由平行文句、同書用語、
傳抄習慣及上下文結構，考察何種
讀法較為可取。但每一項標準都有
適用範圍。例如，短句可能是古老的
形式，也可能是抄手省略的結果；
較通順的語句可能經過後人潤飾，
也可能恰好保留了完整表述。

有效的論證還應容納可能的反證。
假如發現一件年代可靠、流傳系統
獨立的新材料，且它也保存乙本的
較長語句，那麼目前支持甲本的
理由便需要重新檢討。反之，若
找到較早的旁證，而它與甲本
的短句一致，也會改變目前的
證據分布。編者不必因存在不確定性
便拒絕判斷，但應將結論限制
在證據所能支持的程度之內。

回到寫本本身同樣重要。字形是否
完整、墨跡是否連續、是否有刮削
改寫的痕跡，常常比抽象地討論
傳抄規律更加直接。良好的校勘
研究需要在文本結構與物質證據
之間往返，而不是將兩者割裂。

\section{校勘記的作用}
如果全篇異文不多，可以在正文
附近以腳註或表格記錄；當傳本
眾多、異文密集時，則可能需要
專門的校勘記。無論採取何種形式，
主要原則都是使讀者清楚辨認
各寫本所存的文字，以及編者
據以作出選擇的理由。

有些研究目的要求並列不同版本，
而不是勉強合成唯一通行文本。
對寫本流傳史的研究而言，
異文之間的差異本身就是研究對象。
排版形式應服從學術問題，
而不應反過來決定研究者
能夠提出甚麼問題。

\chapter{校訂文與學術表達}\label{chap:edition}
\section{一項暫定選擇}
在這個假設案例中，我們暫以甲本
「謹守其書」作為通讀文，並於
校記中交代乙本增字、丙本闕文。
這項選擇並不表示甲本具有
無條件的優先地位，而是說在
目前材料與研究目的之下，短句
可以構成一種可供檢驗的工作方案。

一份可靠的校訂本，並不要求每項
判斷都具有絕對的確定性。更重要的
是，後來的讀者可以根據編者留下的
材料記錄與推理說明，重建校訂的過程，
甚至提出不同的結論。寫本名稱、
簡號、頁碼、圖版與參考文獻應盡可能
準確，避免因排版整潔而掩蓋資料不足。

\section{章節層級與註釋}
書稿的結構有助於呈現思想的層次。
第~\ref{chap:problem}~章先說明
研究對象，第~\ref{chap:variants}~章
集中比較異文，本章則討論校訂及
其呈現方式。這種交叉引用使讀者
在閱讀後文時，仍能方便地回到
前面的材料說明。

章首頁、頁眉及目錄也承擔導航功能。
它們不是單純的裝飾，而是讓讀者
理解一本書如何組織材料的基本工具。
腳註可以補充來源、解釋術語或指出例外，
卻不宜承擔全部的主要論證。\footnote{可將
此註替換為真實的材料出處、前人意見
或篇幅較短的文字校勘。}

若實際研究涉及表格、圖片、古文字
或不同方向的書寫，可逐項啟用 IMPE
的相應功能。對於不需要的功能，
維持原生 \LaTeX{} 的簡單結構，
通常更利於長期維護。範例不應
迫使使用者接受不必要的設定。

\section{從範例到正式書稿}
使用這份模板時，可以先替換書名、
作者及各章內容，保留適合自己研究
的章節層級。前言之後可增設凡例、
圖版說明、縮略語表，正文之後可
補入參考書目及索引。編輯過程中
應定期檢查 PDF 的分頁、字形、
章節間距及註釋位置，避免到最後
才發現版面與實際內容不相適應。

最值得保留的不是示例文字，而是
清晰的寫作流程：由材料提出問題，
以可檢查的證據進行分析，最後
明示研究結論與尚待解決的疑點。
這一原則既適用於文獻校勘，
也適用於其他長篇學術書稿。

\backmatter
\chapter*{結語}
本書以虛構材料串聯文獻描述、
異文比勘及校訂選擇，展示
一份小型學術著作可以如何組織。
明確的文字層級與舒適的頁面
並不能取代考證，但它們能夠
使論證更易閱讀、檢索與檢驗。

這份原始碼僅是起點。真正的書稿
應增加經過核實的研究材料、
完整的前人研究回顧，以及符合
出版需求的圖表、註釋和書目；
排版則應隨着內容發展而調整。
\end{document}
